\documentclass[12pt]{article}
\usepackage[utf8]{inputenc}
\usepackage[T1]{fontenc}
\usepackage{lmodern}
\usepackage[left=1in,right=1in,top=1.5in,bottom=1.5in]{geometry}
\usepackage{setspace}
\onehalfspacing
\usepackage{amsmath,amssymb}
\usepackage{hyperref}

\begin{document}
\thispagestyle{empty}

\begin{center}
    {\LARGE \textbf{Auditable Memory and Financial Identification Limits in Electronic Limit Order Books}}\\[2.5em]
    
    \textit{[Blinded for Peer Review]}\\[2.0em]
    
    \today\\[3.0em]
\end{center}

\begin{abstract}
\noindent We establish the fundamental financial and econometric limits of identifying hidden adaptive memory in electronic limit order books. We prove that public ticker feeds cannot distinguish competing algorithmic cohorts sharing aggregate sensitivity over finite horizons, and that separating two memory rates differing by $\delta$ requires active executed quantity of order $\Theta(\delta^{-2})$ or financial capital of order $\Theta(\delta^{-4})$. Passive observation cannot substitute for active probing in the absence of new exogenous excitation. Using spectral and finite-horizon Rayleigh--H\"older duality, we certify a $61.44$~USDT infinite-horizon capital lower bound at $\delta=0.05$ and close the finite-horizon non-convex fee and margin duality gap to $0.98\%$ at horizon $L=80$. For unknown two-cohort mixtures, an exact moment-matched Le Cam two-point construction establishes a global non-asymptotic $\Omega((\Delta\eta)^{-6})$ minimax sample complexity lower bound across all measurable estimators, paired with a matching $\mathcal{O}((\Delta\eta)^{-6})$ cubic moment-inversion upper bound on known-midpoint families. We validate these scaling laws across genuinely trained closed-loop adaptive agents (Recursive Least Squares, multi-layer perceptrons, profit-gradient learners, and tabular Bellman Q-learners), showing that settled learning dynamics obey merged-measure equivalence to within $10^{-5}$ relative gap. Finally, evaluating an empirical 768-unit panel and a 66-date multi-day panel, we demonstrate through a nested econometric hierarchy that $>99\%$ of the out-of-sample predictive gain from regulatory inventory disclosure is explained by aggregate level persistence rather than memory filtering.
\end{abstract}

\vspace{2.0em}
\noindent \textbf{Keywords:} Algorithmic trading, system identification, market microstructure, minimax sample complexity, active probing, auditable memory, limit order books. \\[0.8em]
\noindent \textbf{JEL Classification Codes:} C32, C58, D83, G10, G14.

\end{document}
